\documentclass[10pt,a4paper]{article}
\usepackage[margin=2cm]{geometry}
\usepackage[T1]{fontenc}
\usepackage[utf8]{inputenc}
\usepackage{lmodern}
\usepackage{microtype}
\usepackage{amssymb}
\usepackage{booktabs}
\usepackage{tabularx}
\usepackage{array}
\usepackage{multirow}
\usepackage[table]{xcolor}
\usepackage{graphicx}
\usepackage{pgfplots}
\pgfplotsset{compat=1.18}
\usepgfplotslibrary{groupplots}
\usepackage{enumitem}
\usepackage{caption}
\usepackage{subcaption}
\usepackage{float}
\usepackage[hidelinks]{hyperref}
\usepackage{url}

\setlist{nosep,leftmargin=1.4em}
\captionsetup{font=small,labelfont=bf}
\setlength{\parskip}{3pt}
\setlength{\parindent}{0pt}

% ---------------------------------------------------------------- submission details
% Fill in before submitting (Canvas group number; video link, which the course requires in the report).
\newcommand{\groupnumber}{\textcolor{red}{[group number]}}
\newcommand{\videolink}{\textcolor{red}{[presentation video link]}}
\newcommand{\repourl}{\url{https://github.com/X-Bulow/seethrough}}

% ---------------------------------------------------------------- detection-matrix cells
\definecolor{cellD}{HTML}{D6EEDA}
\definecolor{cellM}{HTML}{F6D6D1}
\definecolor{cellU}{HTML}{E4E4E0}
\definecolor{seriesA}{HTML}{2A78D6}
\definecolor{seriesB}{HTML}{EB6834}
\newcommand{\cD}{\cellcolor{cellD}D}
\newcommand{\cM}{\cellcolor{cellM}M}
\newcommand{\cMs}{\cellcolor{cellM}M$^{*}$}
\newcommand{\cU}{\cellcolor{cellU}U}
\newcommand{\cNA}{--}
\newcommand{\rot}[1]{\rotatebox{90}{#1}}
\newcommand{\code}[1]{\texttt{#1}}

\title{\textbf{Do Memory-Safety Detectors See Through Custom Allocators?}\\[2pt]
\large A Detection and Overhead Comparison of ASan, Valgrind Memcheck and HWASan\\[4pt]
\normalsize CS5439 Assignment 2 --- Final Report (Group \groupnumber)}
\author{Cheng Zefeng \quad Wei Yuning \quad Xu Ruofan\\
\small\texttt{bulow@u.nus.edu} \quad \texttt{weiyuning@u.nus.edu} \quad \texttt{e1536075@u.nus.edu}}
\date{November 2026}

\begin{document}
\maketitle
\vspace{-1.5em}
\begin{center}\small
Presentation video: \videolink \qquad Code, corpus and raw results: \repourl
\end{center}

\begin{abstract}
\noindent
Memory-safety detectors such as AddressSanitizer (ASan) and Valgrind Memcheck check every access
against the block that \code{malloc} returned. Many C programs instead take one large block from
\code{malloc} and carve objects out of it themselves. We measure what that costs the detectors.
We built a corpus of 16 heap bugs (overflow, underflow, use-after-free, invalid and double free),
each with a negative control that differs in one line, and ran every bug against seven allocator
layers: \code{malloc}, two allocators of our own (a bump arena and a size-class free list), two of
SQLite's allocators (memsys5 and lookaside), and libc rewrites of our two allocators. Each custom
layer was run with and without the tool's annotation interface. Across 448 bug images and 448 controls
there were no false positives. Without annotation, every detector found only 1--2 of 14--16 bugs
inside a custom allocator, against 12--15 of 13--15 on \code{malloc}. Annotation recovered most of
them: 10--11 for ASan, 13--14 for Valgrind and 9--11 for HWASan. Each annotation interface has its
own ceiling: poisoning cannot catch an overflow into another live object or a stale pointer to
reissued storage, and neither poisoning nor tagging can express a bad \code{free}, while memory
tagging is only 16-byte precise. The ASan and Valgrind verdicts are identical, cell for cell, on a
second machine. On SQLite and Lua benchmarks ASan costs 2.2--3.0$\times$ in time and Valgrind
24--29$\times$; switching SQLite to its memsys5 allocator lowers ASan's memory overhead from
7.9$\times$ to 1.4$\times$, which is the other side of the visibility it loses.
\end{abstract}

\input{src/01-introduction}
\input{src/02-background}
\input{src/03-related-work}
\input{src/04-methodology}
\input{src/05-detection}
\input{src/06-performance}
\input{src/07-discussion}
\input{src/08-conclusion}

\bibliographystyle{plainurl}
\bibliography{src/refs}

\appendix
\input{src/appendix}

\end{document}
